\begin{afterword}
\summaryhead

The post introduces entropy. A system with eight equally likely states has an entropy of 3 bits: three yes-or-no questions find its state. A system Y whose four states have probabilities 1/2, 1/4, 1/8 and 1/8 has an entropy of 1.75 bits. A code with a one-symbol word for the likeliest state and three-symbol words for the rarest needs 1.75 questions on average. The general formula is the sum of $-p \log_2 p$ over the states. Giving Y4 the shorter word 10 would make messages ambiguous, so ``short words are a conserved resource.'' In the limit, a message's length corresponds to its probability; this is the Minimum Description Length or Minimum Message Length form of Occam's razor.

So, the author argues, the labels we give concepts are not quite arbitrary, and ``You can X any way you like'' obstructs learning to X wisely. Basic-level categories such as ``chair'' are the level people talk at, and they have shorter names than ``recliner'' or ``furniture'': frequent use goes with short words.

\responsehead

The first two-thirds of the post are a short and correct introduction to entropy and codes. Every number is correct: the code for X, the 1.75 bits for Y, the 0.375 contributed by a state of probability 1/8, and the two readings of ``1001'' under the altered code. The post also states the limits it needs: a perfect code is not always possible, and coding many copies together comes arbitrarily close. Its moral, that short words are a conserved resource, is an informal statement of the Kraft--McMillan inequality.

The last third applies this to language. The post's own evidence there is one example, chair against recliner and furniture, but the literature supports the claim. It is Zipf's law of abbreviation, confirmed in almost a thousand languages, and a 2011 study found that word length follows a word's information content even more closely than its frequency. The claim is better supported than the post shows.

One step needs qualifying. The route through basic-level categories is less direct than the post suggests. Rosch and colleagues found basic names preferred when people name objects, but in nine of fifteen comparisons the more general word was the more frequent one, and their study did not test the claim about length.

A reader should keep apart two probabilities that sit side by side in the last third. How often a word is used sets its length in a language; how probable a hypothesis is is what Minimum Message Length measures. The post does not equate them, and the author separates them explicitly in ``Occam's Razor,'' where ``witch'' is short but names ``extraordinary assertions.''

In short: a correct introduction to entropy and short codes, whose claim about the length of words is sound and better supported than its one example shows, with a route through basic-level categories that is less direct than it suggests.
\end{afterword}
